\chapter{Methodology}
\label{chap:methodology}

This chapter presents the methodology used to design and evaluate the
Personalized Travel Guide Assistant based on Retrieval-Augmented Generation
(RAG). The system consists of two principal workflows: offline tourism
knowledge-base construction and online personalized travel-query processing.

The offline workflow collects tourism webpages and documents, converts them
into a consistent textual format, removes irrelevant content, divides the
documents into retrieval units, enriches the resulting chunks with metadata,
generates embeddings, and stores them in the knowledge base.

At query time, the system analyzes the user's request and conversational
context, retrieves candidate chunks using dense vector search and BM25,
combines the results using hybrid fusion, reranks the strongest candidates,
checks whether the evidence covers the user's requirements, performs bounded
recovery retrieval when necessary, and generates a grounded personalized
response with source attribution.

The methodology treats retrieval as a multi-stage process because final-answer
quality depends on both the relevance and the coverage of the supplied
evidence. It also treats query understanding, retrieval, and answer generation
as separately inspectable stages so that errors can be localized rather than
attributed only to the final language model.

\section{System Architecture}

The proposed system consists of five principal layers:

\begin{enumerate}[label=(\roman*)]
    \item \textbf{Knowledge-base construction}, including tourism data
    acquisition, document conversion, cleaning, chunking, metadata enrichment,
    embedding generation, and database storage;

    \item \textbf{Query understanding}, including request routing,
    conversational query rewriting, intent identification, operation
    identification, explicit-constraint extraction, and retrieval-facet
    extraction;

    \item \textbf{Retrieval and ranking}, including dense vector retrieval,
    BM25 lexical retrieval, Reciprocal Rank Fusion, metadata and optional
    geographic signals, and CrossEncoder reranking;

    \item \textbf{Contextual processing}, including persistent user memory,
    temporary conversation state, retrieval planning, evidence checking,
    recovery retrieval, and missing-evidence diagnosis; and

    \item \textbf{Grounded response generation}, including context selection,
    answer-readiness checking, safety and constraint handling, Markdown
    generation, source attribution, and knowledge-gap handling.
\end{enumerate}

These layers are implemented as separate components so that each stage can be
inspected, modified, and evaluated independently.

\section{Tourism Knowledge-Base Construction}

\subsection{Tourism Data Acquisition}

The first stage constructs a tourism-specific knowledge base from supported
websites and documents. Tourism information is collected from sources such as
official tourism portals, destination websites, online travel guides, and PDF
documents.

Website-specific crawlers collect relevant tourism URLs while applying URL
normalization and duplicate controls. The crawlers restrict collection to
relevant sections of the supported websites rather than following all
available links indiscriminately. Each collected URL retains its source and
processing status so that ingestion progress can be monitored.

For document-based sources, the system supports PDF-to-Markdown conversion.
The conversion attempts to preserve useful structural information, including
headings, paragraphs, and lists. This structure is later used during document
chunking.

Each source document is stored with information including:

\begin{itemize}
    \item document type and source location;
    \item extraction method;
    \item raw and cleaned Markdown;
    \item ingestion status;
    \item chunking method;
    \item embedding model; and
    \item creation and update timestamps.
\end{itemize}

Local files are identified using file hashes, while web documents are
identified using their normalized source URLs. These identifiers reduce
duplicate ingestion.

\subsection{Content Extraction and Cleaning}

HTML pages are converted into Markdown before ingestion. The extraction
process focuses on the main page content while attempting to preserve useful
headings, lists, links, and tables.

Raw tourism content may contain elements that are not useful for retrieval,
including navigation menus, advertisements, cookie notices, social-media
controls, repeated headers and footers, unrelated links, duplicated text, and
empty or low-information sections. The cleaning stage removes these elements,
normalizes whitespace and line breaks, and preserves the structural information
needed for chunking. Irrelevant image markup is removed from the textual
knowledge base. When useful OCR text is available from document images, it may
be retained as textual context.

Documents or sections containing insufficient useful information are excluded
before chunk storage.

\section{Structure-Aware Document Chunking}

Complete tourism articles are generally too large and topically diverse to be
used directly as retrieval units. The cleaned documents are therefore divided
into smaller chunks.

The implemented method first divides Markdown according to heading levels from
level one to level four. Each section is then processed using recursive
character-based splitting. The splitter prioritizes boundaries such as
paragraphs, line breaks, sentences, punctuation, and spaces.

The principal chunking configuration is shown in
Table~\ref{tab:chunk-configuration}.

\begin{table}[H]
\centering
\caption{Document chunking configuration}
\label{tab:chunk-configuration}
\begin{tabular}{ll}
\hline
\textbf{Parameter} & \textbf{Value} \\
\hline
Maximum chunk size & 1,200 characters \\
Chunk overlap & 150 characters \\
Minimum candidate length & 150 characters \\
Minimum accepted word count & More than 18 words \\
Heading context & Included \\
\hline
\end{tabular}
\end{table}

Short adjacent chunks from the same section can be merged when the combined
text remains within the configured maximum size. Chunks containing insufficient
information are removed.

For a document organized as:

\begin{verbatim}
# Hanoi
## Hoan Kiem Lake
### Things to do
\end{verbatim}

the relevant heading context is retained with the chunk. This prevents a
passage describing an activity from losing its relationship with the
attraction or destination. Each accepted chunk is assigned a normalized hash
to support duplicate detection and stable identification.

\section{Tourism Metadata Enrichment}

Each chunk is enriched with structured tourism metadata before being stored.
Metadata extraction is performed in batches using an LLM-assisted extraction
process. Context from the preceding chunk, such as its heading, summary,
country, and city, can be supplied to improve continuity across related
sections.

The extracted metadata includes:

\begin{itemize}
    \item country, region, city, and province;
    \item place name and place type;
    \item semantic tags;
    \item activities and travel styles;
    \item suitable traveler groups;
    \item chunk topic and summary;
    \item extraction confidence and reasoning; and
    \item source and update information when available.
\end{itemize}

Controlled vocabularies and normalization rules are applied to categorical
values such as place type, activity, travel style, suitable traveler group,
query intent, operation, and budget level. Controlled values improve
consistency, but they are not assumed to represent every possible phrase a user
may enter. The LLM parser therefore maps natural-language expressions to the
closest supported value and retains explicit free-text constraints when a
controlled value would lose important meaning.

\subsection{Geographic Enrichment}

Geographic coordinates can be added after metadata extraction. Chunks are
grouped using a normalized combination of place name, city, province, and
country. Each unique place is geocoded once, and the resulting coordinates are
assigned to other chunks associated with the same place.

The geocoding request includes available city, province, and country
information to reduce ambiguous matches. When multiple candidates are
returned, they are scored using place-name and location agreement. Coordinates
are stored only when a suitable result is available. This process avoids
requesting coordinates independently for every chunk and reduces repeated
geocoding operations.

\section{Embedding Generation and Knowledge Storage}

Each enriched chunk is transformed into a dense vector using the
\texttt{all-MiniLM-L6-v2} sentence-embedding model. The model produces a
384-dimensional representation. The same embedding model is used for both
document chunks and user queries.

Dense retrieval compares the query and chunk embeddings using cosine
similarity:

\begin{equation}
\operatorname{sim}(Q,c_i)
=
\frac{\mathbf{v}_Q \cdot \mathbf{v}_i}
{\lVert\mathbf{v}_Q\rVert\lVert\mathbf{v}_i\rVert},
\end{equation}

where $\mathbf{v}_Q$ is the query embedding and $\mathbf{v}_i$ is the
embedding of chunk $c_i$.

The knowledge base is implemented using PostgreSQL with vector-search support.
Each stored chunk contains its text and hash, document identifier, section
heading, vector embedding, tourism metadata, optional geographic coordinates,
source-document relationship, and creation and update information. Relational
fields support metadata filtering, while the vector field supports
nearest-neighbor retrieval.

\section{Request Routing and Query Understanding}

Before retrieval, the system determines whether the user's request belongs to
the tourism domain. Requests are classified into three categories:

\begin{itemize}
    \item \textbf{travel}, when the retrieval pipeline should continue;
    \item \textbf{ambiguous}, when the system should request clarification; or
    \item \textbf{out of scope}, when the request is unrelated to travel
    assistance.
\end{itemize}

If the routing model is unavailable, the system continues with guarded
retrieval. The later evidence-checking stage still prevents unsupported factual
answers.

\subsection{Conversational Query Rewriting}

Travel conversations commonly contain follow-up questions such as ``What about
there?'', ``Which one is quieter?'', or ``Where should I stay for the final two
nights?''. These questions depend on previous messages and may not be suitable
for direct retrieval.

The query-rewriting stage uses recent conversation history to transform the
current message into a standalone retrieval query. It resolves references while
preserving the user's intent, destinations, dates, activities, budgets, and
other explicit constraints. The rewriting stage is instructed not to copy
unrelated constraints from previous messages and not to introduce preferences
or facts absent from the conversation.

\subsection{Query Parsing}

The standalone query is parsed into a structured representation containing:

\begin{itemize}
    \item query intent and requested operation;
    \item explicit constraints and negative constraints;
    \item country, region, city, and province;
    \item place types, activities, and travel styles;
    \item suitable traveler groups;
    \item budget, trip duration, and dates;
    \item nearby-place requirements; and
    \item maximum-distance requirements.
\end{itemize}

Intent identifies the travel domain of the request, such as accommodation,
food, transport, or itinerary planning. Operation identifies what the system
must do with the request, such as lookup, recommendation, comparison, planning,
or explanation. Keeping these two concepts separate avoids treating, for
example, every accommodation question as the same type of operation.

\subsection{Explicit Constraints and Retrieval Facets}

The parser separates \textit{explicit constraints} from \textit{retrieval
facets}. Explicit constraints represent requirements directly stated by the
user, such as a named destination, trip duration, budget, dietary restriction,
or an activity that must be avoided. They are retained for understanding
evaluation and final-answer compliance.

Retrieval facets are normalized semantic signals, including place type,
activity, travel style, and suitable traveler group. These facets help retrieve
and rank useful evidence, but most are not treated as strict database filters.
The separation prevents a soft preference such as an interest in photography
from eliminating otherwise relevant evidence.

Positive and negative constraints are also distinguished. A phrase such as
``avoid party nightlife'' must not be converted into a positive nightlife
facet. Negative constraints are preserved for evidence selection and response
generation. Durable avoidance preferences may also be retained in persistent
memory, whereas trip-specific restrictions remain in conversation state.

Known city names and controlled tourism vocabularies validate extracted values.
Explicitly named destinations are additionally identified through deterministic
text matching. This hybrid LLM and rule-based normalization reduces the risk
that the parser omits a location stated directly in the query or returns a
structurally invalid constraint.

\subsection{Region and Location Normalization}

Broad areas such as Northern Vietnam, Central Vietnam, and Southern Vietnam are
represented as regions rather than provinces. Common aliases are normalized to
canonical region names. When only a broad region is specified, the parser must
not invent a city or store the region as a province.

A region is not applied directly as an exact province filter because it may
contain several provinces and cities and because corpus metadata may represent
the same area at different geographic levels. Instead, it is preserved as
query context and may be used to construct destination-scoped retrieval tasks.
When a user later identifies a specific city, that city becomes the primary
location constraint and an incompatible broad province value is removed.

\section{User Memory and Conversation State}

The system distinguishes between persistent user memory and temporary
conversation state.

\subsection{Persistent User Memory}

Persistent memory stores selected long-term preferences and user
characteristics, including preferred travel styles, preferred activities,
interests, budget preference, conditions to avoid, preferred answer length,
preferred tone and explanation style, expertise, and selected non-sensitive
personal facts.

Memory items are stored with a user identifier, memory type, normalized
content, importance value, activity status, and timestamps. Active memories are
ordered by importance and recency before being converted into a structured
user profile. Travel-style and activity preferences can provide small ranking
signals. The complete relevant profile can also be used by answer generation.

Not every user statement is stored permanently. Memory extraction is
restricted to information likely to remain useful across future conversations.

\subsection{Temporary Conversation State}

Conversation state stores trip-specific information relevant only to the
current discussion. It includes a bounded conversation summary, current
destination, travel dates, trip duration, temporary budget, selected places,
and trip-specific constraints.

Separating temporary state from persistent memory prevents a short-term trip
requirement from being interpreted as a permanent preference. State size is
bounded to reduce prompt growth. If an LLM state update is truncated or
malformed, the previous valid state is retained instead of replacing it with
unvalidated content.

\section{Retrieval Planning}

Simple factual questions can normally be processed using one retrieval query.
More complex requests may contain several independent information requirements.
The retrieval planner analyzes the standalone query and creates focused
retrieval tasks when decomposition is useful. For example, a comparison between
two destinations may produce one destination-specific search for each city. An
itinerary request may require evidence about attractions, activities, food, and
practical travel information.

The planner distinguishes evidence requirements from response-format
instructions. A request for a concise answer or a three-day structure should
guide answer construction but should not automatically become a separate
retrieval requirement unless factual evidence is needed.

The implementation limits the planner to a maximum of six evidence requirements
and six retrieval tasks. If planning fails, the system falls back to a general
search or destination-scoped searches for comparison queries. This bounded
approach provides decomposition without allowing an excessive number of
retrieval operations.

\section{Metadata-Aware Retrieval}

The structured query representation is converted into retrieval filters and
ranking signals. Supported exact filters include country, city, province, and,
when sufficiently reliable, place type.

Hard place-type filtering is used only when the intent clearly requires a
particular category, such as attraction, accommodation, food, or event search.
For broader requests, place types, activities, travel styles, suitability
information, and regions are used as retrieval context or ranking signals
rather than strict filters.

If optional province or place-type filters return no candidates, the system
removes these optional restrictions and repeats retrieval. Every relaxation is
recorded in the pipeline trace. This process prevents complete retrieval
failure when metadata is incomplete or the query has been interpreted too
narrowly.

\section{Dense Vector Retrieval}

Dense retrieval compares the query embedding with stored chunk embeddings. It
is useful when the query and relevant documents express the same concept using
different words. For each retrieval task, the dense retriever returns the most
semantically similar chunks from the metadata-filtered collection. The
pipeline-level ceiling is 30 dense candidates, although the retrieval planner
normally supplies a smaller task-specific limit.

\section{BM25 Lexical Retrieval}

BM25 provides lexical matching between the query and stored tourism chunks. The
searchable representation combines place name, summary, and original chunk
text. BM25 is particularly useful for a specific attraction, landmark,
district, or uncommon local term. The pipeline-level ceiling is 30 BM25
candidates, subject to the smaller limit assigned to the current retrieval
task.

\section{Hybrid Retrieval}

Dense and BM25 retrieval have complementary strengths. Their result lists are
combined using Reciprocal Rank Fusion (RRF). For candidate chunk $c$, the
fusion score is:

\begin{equation}
S_{\mathrm{RRF}}(c)
=
\frac{\mathbb{I}_{\mathrm{dense}}(c)}{k+r_{\mathrm{dense}}(c)}
+
\frac{\mathbb{I}_{\mathrm{BM25}}(c)}{k+r_{\mathrm{BM25}}(c)},
\end{equation}

where $r_{\mathrm{dense}}(c)$ and $r_{\mathrm{BM25}}(c)$ are the ranks of $c$
in the two lists. The implementation uses $k=60$. Candidate chunks are
deduplicated using stable chunk identifiers before ordering.

Small additional ranking signals can be applied for matching query activities,
travel styles, suitable traveler groups, stored user preferences, geographic
proximity, and source freshness. These signals are intentionally smaller than
the main retrieval contribution so that they refine rather than replace textual
relevance.

\subsection{Location-Aware Retrieval}

Location relevance is primarily handled through structured country, city, and
province metadata. When the query explicitly names a city, unrelated cities
can be excluded from the initial search. Broad regions remain contextual rather
than being converted into incorrect exact province filters.

The knowledge base can additionally store latitude and longitude values. The
retrieval pipeline supports a small distance-based ranking adjustment when both
target and candidate coordinates are explicitly available. Coordinate-based
ranking is an optional supporting signal rather than a principal retrieval
mechanism.

\section{CrossEncoder Reranking}

Hybrid retrieval improves candidate coverage, but its order may still contain
passages that are only partially relevant. The system therefore applies a
CrossEncoder reranking stage using
\texttt{cross-encoder/ms-marco-MiniLM-L-6-v2}. Unlike the embedding model, the
CrossEncoder processes query and passage jointly.

Each planned search normally requests between 5 and 20 candidates from each
first-stage retriever, with a fallback value of 10. After results from all
planned searches have been merged, up to 20 fused candidates are forwarded to
the CrossEncoder. The reranker retains up to eight chunks as final evidence.

Destination-comparison queries receive additional balancing. The shortlist and
final evidence selection attempt to include evidence from each explicitly
compared city. Balancing affects candidate membership, while CrossEncoder
scores determine relevance order. This procedure cannot create evidence absent
from the corpus; such cases are handled by coverage checking and corpus-gap
diagnostics.

\begin{table}[H]
\centering
\caption{Configuration of the active retrieval and reranking pipeline}
\label{tab:retrieval-configuration}
\begin{tabular}{p{0.37\textwidth} p{0.52\textwidth}}
\hline
\textbf{Component} & \textbf{Configuration} \\
\hline
Embedding model & \texttt{all-MiniLM-L6-v2} \\
Embedding dimension & 384 \\
Candidates per planned search & 5--20 per dense and BM25 search; fallback 10 \\
Pipeline retrieval ceiling & Up to 30 candidates per retriever \\
Hybrid fusion & RRF with $k=60$ and small metadata signals \\
CrossEncoder input & Up to 20 fused candidates \\
Reranking model & MS MARCO MiniLM-L6-v2 CrossEncoder \\
Final evidence & Up to 8 reranked chunks \\
Comparison handling & Candidate and evidence balancing across compared cities \\
\hline
\end{tabular}
\end{table}

\section{Evidence Coverage and Recovery Retrieval}

A high relevance score does not guarantee that all parts of a complex request
are supported. After reranking, the system checks selected evidence against the
requirements produced by the planner. Each requirement is classified as:

\begin{itemize}
    \item \textbf{covered}, when the evidence is sufficient;
    \item \textbf{partially covered}, when only part is supported; or
    \item \textbf{not covered}, when no directly useful selected evidence is
    available.
\end{itemize}

Supporting chunk identifiers must be provided for covered or partially covered
requirements. This prevents the checker from declaring coverage without
identifying corresponding evidence. Coverage ratio assigns full credit to a
covered requirement, partial credit to a partially covered requirement, and no
credit to a not-covered requirement.

When evidence is missing, the checker can generate focused recovery queries.
The system executes at most three recovery searches. New candidates are merged
with original results, followed by fusion, reranking, and a second evidence
check. The trace records whether recovery added unique candidates, changed the
selected evidence, or improved requirement coverage.

\subsection{Missing-Evidence and Corpus Diagnostics}

Failure to select evidence does not automatically prove that the corpus lacks
the requested information. The system therefore records a diagnostic for every
unresolved requirement after recovery:

\begin{itemize}
    \item \texttt{not\_probed}, when no focused recovery query was executed and
    no corpus-gap conclusion can be made;
    \item \texttt{likely\_corpus\_gap}, when the focused probe returns no
    candidate;
    \item \texttt{likely\_corpus\_content\_gap}, when candidates exist but do
    not contain the required information;
    \item \texttt{likely\_reranking\_failure}, when the focused probe finds new
    candidates that are not selected; or
    \item \texttt{selected\_evidence\_insufficient}, when probe evidence is
    selected but remains insufficient.
\end{itemize}

These labels are diagnostics, not evidence-coverage grades. In particular,
\texttt{not\_probed} must not be interpreted as proof of missing corpus data.
Response-format requirements and personalization facets are excluded from
evidence coverage unless they require an external factual claim.

\section{Answer Readiness and Knowledge-Gap Handling}

After retrieval, the system assigns one of three answer-readiness modes:

\begin{itemize}
    \item \textbf{complete}, when the evidence sufficiently covers the factual
    requirements;
    \item \textbf{partial}, when useful evidence exists but one or more
    requirements remain incompletely supported; or
    \item \textbf{insufficient}, when no adequate evidence supports normal
    answer generation.
\end{itemize}

In partial mode, the answer generator must separate supported information from
missing information. In insufficient mode, normal generation is blocked and
the system explains that the current database does not contain enough verified
information. The unresolved topic is recorded as a knowledge gap. This gate
reduces the likelihood that the LLM replaces failed retrieval with unsupported
internal knowledge.

\section{Grounded Response Generation}

The final generation context includes the original query, rewritten standalone
query, parsed constraints and facets, selected evidence, relevant persistent
memory, temporary conversation state, evidence coverage, missing-evidence
diagnostics, and answer readiness. Each evidence chunk receives an evidence
identifier and retains its place name, source location, and relevant metadata.

The answer model is instructed to:

\begin{itemize}
    \item base factual travel claims on selected evidence;
    \item apply user preferences only when relevant;
    \item respect explicit positive and negative constraints;
    \item avoid inventing unsupported details;
    \item identify missing or uncertain information;
    \item avoid exposing retrieval scores or internal pipeline details; and
    \item produce a readable Markdown response.
\end{itemize}

For safety-sensitive food constraints, including allergies, shellfish
avoidance, ingredient uncertainty, and cross-contact risk, the model must not
claim that an option is safe unless the evidence verifies the relevant facts.
Otherwise, it must state the limitation and recommend direct confirmation with
the provider.

\section{Source Attribution}

Every tourism chunk maintains a relationship with its original document and
source location. Sources associated with selected evidence are included in the
response payload and duplicate sources are consolidated before presentation.
Source attribution improves transparency, but the presence of a source does
not by itself guarantee every generated statement is correct. It is therefore
combined with evidence-constrained prompting and coverage checking.

\section{Retrieval Debugging and Telemetry}

A retrieval-debugging interface exposes intermediate pipeline results rather
than displaying only the final answer. It supports inspection of rewritten
queries, parsed constraints and facets, retrieval plans, dense and BM25
candidates, fusion results, applied filters and relaxations, metadata signals,
CrossEncoder scores, coverage assessments, recovery searches, corpus-gap
diagnostics, selected evidence, readiness, and the generated answer.

Operational telemetry is recorded for reproducibility and efficiency analysis.
Depending on model-provider reporting, each stage stores:

\begin{itemize}
    \item latency and model identifier;
    \item input, cached-input, and output token counts;
    \item total token usage;
    \item estimated cost based on a versioned pricing configuration;
    \item pipeline and trace versions;
    \item generator and judge prompt hashes; and
    \item corpus fingerprint and relevant model configuration.
\end{itemize}

Estimated cost is an experimental estimate rather than a billing record. A
corpus fingerprint connects results to the exact knowledge-base snapshot used
for a run.

The exact model assignments, dataset construction process, evaluation metrics,
judge configuration, and reproducibility controls are specified in the
experimental-setup chapter. Keeping these choices separate from the proposed
method makes it possible to change an experimental model without redefining the
architecture itself.

\section{End-to-End Query Processing}

The complete online process is summarized as follows:

\begin{enumerate}
    \item The request router classifies the request as travel, ambiguous, or out
    of scope.
    \item Persistent user memory and temporary conversation state are loaded.
    \item The current message is rewritten into a standalone query.
    \item The parser extracts intent, operation, location, explicit constraints,
    negative constraints, and retrieval facets.
    \item Region aliases and deterministic destination matches are normalized.
    \item The retrieval planner creates bounded evidence requirements and search
    tasks.
    \item Applicable hard filters and semantic ranking signals are constructed.
    \item Dense and BM25 retrieval are executed for each task.
    \item Optional filters are relaxed if they produce no candidates.
    \item Candidates are deduplicated and combined using RRF and small metadata
    signals.
    \item The strongest fused candidates are CrossEncoder-reranked.
    \item The evidence checker evaluates requirement coverage.
    \item Up to three recovery searches are performed when evidence is missing,
    followed by fusion, reranking, and rechecking.
    \item Missing-evidence and corpus diagnostics are recorded.
    \item The evidence gate assigns complete, partial, or insufficient
    readiness.
    \item The LLM generates a personalized grounded response when permitted.
    \item Sources, trace information, latency, tokens, and estimated cost are
    consolidated for output and evaluation.
\end{enumerate}

\section{Methodological Summary}

The methodology treats the Personalized Travel Guide Assistant as a complete,
inspectable retrieval and generation system rather than an isolated LLM
application. The knowledge-construction workflow transforms heterogeneous
tourism sources into cleaned, structure-aware, metadata-enriched, and
vectorized chunks. At query time, the system performs routing, contextual
rewriting, structured parsing, bounded planning, metadata-aware dense and BM25
retrieval, Reciprocal Rank Fusion, CrossEncoder reranking, evidence checking,
recovery retrieval, corpus-gap diagnosis, and readiness-gated answer
generation.

Persistent preferences and temporary state support personalization without
treating every interaction as permanent memory. The evaluation protocol
separates query understanding, retrieval ranking, and final-answer quality,
uses an independent LLM judge with human confirmation, and records sufficient
telemetry to connect results to the exact prompts, models, pipeline version,
and corpus snapshot used in each run.
